%!TEX root = ../manuscript.tex

\begin{theorem}[Certifying bigon-barrier decision]\label{thm:certifying-bigon-decision}
Let $w$ be a signed Gauss word with $n$ chords, and let $\mathcal{B}$
denote the rewriting system of bigon R1-down and bigon R2-down moves.
There is a deterministic algorithm (any maximal greedy $\mathcal{B}$-run)
that, in at most $n$ steps, outputs exactly one of the following two
certificates:
\begin{enumerate}
\item A deletion sequence
  $w = w_0 \xrightarrow{\mu_0} w_1 \xrightarrow{\mu_1} \cdots
  \xrightarrow{\mu_{t-1}} w_t = \emptyset$ of bigon moves,
  certifying that $w$ is bigon-reducible.
\item A nonempty word $z = w_t$ that admits no bigon move
  (equivalently, every directly-nested pair in the nesting forest
  $N(z)$ fails the bigon outer-arc condition, and no adjacent
  opposite-sign pair is present; see Theorem~\ref{thm:stuck-bigon}),
  together with the deletion sequence from $w$ to $z$, certifying
  that $w$ is \emph{not} bigon-reducible.
\end{enumerate}
Both certificates are verifiable in polynomial time. The algorithm
runs in $O(n^4)$ time with a naive move scanner, or $O(n^3)$ time
with incremental maintenance of the available-move set.
\end{theorem}

\begin{proof}
Starting from $w_0 = w$, the algorithm iterates: at each state $w_t$,
if any bigon move is available, it chooses one ($\mu_t$) and sets
$w_{t+1} = \mu_t(w_t)$; otherwise, it halts and outputs $w_t$.

\emph{Termination.} Each bigon move decreases the chord count by at
least one, so the procedure halts after at most $n$ steps.

\emph{Yes-certificate.} If the procedure halts with $w_t = \emptyset$,
the recorded sequence of moves is, by construction, a valid
$\mathcal{B}$-reduction of $w$ to the empty word. This certifies
bigon-reducibility.

\emph{No-certificate.} Suppose the procedure halts with
$z = w_t \neq \emptyset$. By construction, $z$ is stuck (no bigon
move applies at $z$). We show $w$ is not bigon-reducible.

Assume for contradiction that $w$ admits a bigon reduction to
$\emptyset$. By induction on $t$, each intermediate state $w_i$
admits a bigon reduction to $\emptyset$: the base case $i = 0$ is the
hypothesis, and for the induction step, Lemma~\ref{lem:rescue}
(Reduction-Preservation) applied at $w_i$ with move $\mu_i$
shows that $w_{i+1} = \mu_i(w_i)$ also admits a bigon reduction to
$\emptyset$.

Hence $z = w_t$ admits a bigon reduction to $\emptyset$. Since $z$ is
nonempty, any such reduction is nonempty and therefore begins with a
bigon move applied at $z$. But $z$ is stuck, so no bigon move is
available at $z$, yielding a contradiction.

Therefore $w$ is not bigon-reducible, and $(z, (\mu_0, \mu_1, \ldots,
\mu_{t-1}))$ is a valid no-certificate.

\emph{Verification.} A yes-certificate is verified by replaying the
listed moves from $w$ and checking at each step (a) that the listed
move is indeed a valid bigon move at the current word, and (b) that
the final state is $\emptyset$. Each step's validity check runs in
$O(n)$ time for R1-down (adjacency test) or $O(n^2)$ time for
R2-down (direct nesting plus bigon outer-arc emptiness), giving
$O(n^3)$ total verification time. A no-certificate is verified
identically, with the additional step of scanning the residual word
$z$ to confirm that no bigon move applies, at cost $O(n^3)$.

\emph{Running time.} There are at most $n$ deletion steps. A naive
scanner recomputes all candidate moves from scratch at each step,
costing $O(n^3)$ per step for R2-down enumeration (as in
Theorem~\ref{thm:greedy-main}) and giving $O(n^4)$ total. An
incremental data structure maintaining the set of currently-bigon-
eligible pairs and opposite-sign adjacencies can be updated in
$O(n^2)$ amortized time per move, yielding $O(n^3)$ overall.
\end{proof}

\begin{remark}[Why the residual core is a certificate]\label{rem:why-residual-is-cert}
A maximal greedy run in an arbitrary rewrite system typically proves
only local optimality, not global impossibility. What makes
Theorem~\ref{thm:certifying-bigon-decision} stronger is
Lemma~\ref{lem:rescue}: if $w$ were bigon-reducible, every available
first move would preserve bigon-reducibility, and so no greedy run
could terminate at a nonempty stuck word. The stuck residual $z$ is
therefore a global obstruction, not a local one. This is the essential
algorithmic content of Lemma~\ref{lem:rescue} and is what elevates
the stuck-state characterization
(Theorem~\ref{thm:stuck-bigon}) from a structural observation into
part of a complete decision procedure.
\end{remark}

\begin{remark}[Scope: bigon-monotone reducibility, not virtual unknotting]
\label{rem:scope-bigon-vs-unknotting}
We emphasize that Theorem~\ref{thm:certifying-bigon-decision} decides
the specific combinatorial problem of \emph{bigon-monotone
reducibility of signed Gauss words} — whether a word simplifies to
the empty word using only bigon R1-down and bigon R2-down moves. This
is strictly weaker than the virtual unknotting problem: most classical
unknot diagrams are \emph{not} bigon-reducible (for example,
Goeritz and Culprit, as established in the computational study of
Section~\ref{sec:bigon}), yet they are, tautologically, unknots.
Bigon-monotone reducibility is the most restrictive natural
Reidemeister-downward rewrite system on signed Gauss words, and the
exact decidability of this restricted problem is the scope of
Theorem~\ref{thm:certifying-bigon-decision}. Generalizing to the
permissive R2 setting, or to the full Reidemeister rewrite system
with up-moves, raises problems whose complexity status we leave open.
\end{remark}
